\section{Validation}

% Une simulation qui tourne n'est pas une simulation juste.
% Cette section doit donner des NOMBRES, pas des impressions visuelles.

\subsection{Validation du moteur sur un cas de reference}

% La bille dans une boite en rotation. Violation de contrainte mesuree a
% \num{2.2e-16} m, independante de $\dt$. Energie strictement non croissante
% en boite fixe.
%
% Figures disponibles : fig01_respect_contrainte, fig02_energie,
% fig03_entrainement, fig04_trajectoire.
%
% Resultat a mettre en avant : le critere d'entrainement solide n'est verifie
% que dans le regime centrifuge ($\omega^2 R / g \gtrsim 1$), et les
% parametres suggeres par le sujet ne s'y trouvent pas. Bon exemple d'analyse
% critique d'un resultat de simulation.

\subsection{Invariants du modele de foule}

% Non-interpenetration, non-sortie du domaine, conservation du nombre d'agents.
% Donner l'ecart mesure, pas seulement une affirmation.

\subsection{Sensibilite a la discretisation}

% $\dt$ divise par 2 puis par 4 ; $N$ croissant ; plusieurs graines.

\subsection{Comparaison a un resultat connu}

% Diagramme fondamental de Weidmann et debit specifique de goulot de la
% litterature.

\subsection{Coherence dimensionnelle}

% Tableau de toutes les grandeurs avec leurs unites et leurs ordres de grandeur.
